\section{Stage 2：Lasso稳定性选择回答“这个描述符是否依赖少数材料”}

单次筛选在小样本材料数据中非常脆弱。两个 A 族描述符可能都在表达 Na 配位尺度，删除几个样本后 Lasso 可能从平均键长切换到多面体体积；一个空间群代理也可能因某次划分恰好获得高系数。单次“入选”不足以支持结构假设。

Lasso 通过 $L_1$ 惩罚把部分系数压为严格零，因此天然给出“这次子样本中是否选择该描述符”的事件\cite{Tibshirani1996}。框架重复 100 次无放回半样本抽取，在每个子样本内部单独执行中位数插补、标准化和 Lasso($\alpha=0.05$)。稳定性分数是非零系数出现的比例。该思想来自高维统计中的 stability selection，已广泛用于变量和生物标志物筛选\cite{Meinshausen2010}。

这里 Ridge 与 Lasso 的职责必须区分：Ridge 在 Stage 1 中平滑估计体系背景，在 Stage 4 中作为低容量预测探针；Lasso 只在 Stage 2 中制造稀疏选择事件。若用 Ridge 做稳定性选择，几乎所有特征都会得到非零系数，无法定义清晰的入选频率。若只运行一次 Lasso，则相关描述符之间的竞争会使结果依赖一次抽样。

固定加入 15 个标准正态噪声列，回答另一个问题：在 38 个真实描述符和多次搜索中，一个完全无意义的变量最多可以表现多好？真实描述符不仅需要选择频率超过 0.60，还必须超过噪声选择频率的第 95 百分位数。噪声列不是正式假发现率控制，但提供了与当前样本量、候选数和 Lasso 设置一致的经验零基线。

随后按物理族选择代表。这样做不是为了平均分配名额，而是防止平均 Na--X 距离、最大键长和多面体体积把同一局域尺度重复包装成多个“独立发现”。二元模式每族最多一个代表；三元模式每族最多两个，第二个名额仅用于“同族两个互补量 + 相邻族一个量”的受限公式，不被视作第二项独立科学结论。

\section{Stage 3：受约束组合搜索表达多因素迁移，而不制造任意公式}

Na$^+$ 传输通常需要多个条件同时满足。较弱的局域束缚若没有连续 Na 位点网络，可能仍无法形成长程扩散；大量空位若对应高度弯曲或低维通道，也未必提高宏观电导率。因此，组合描述符具有物理必要性。但从 38 个特征任意生成加、减、乘、除、对数和幂函数，会使搜索空间急剧膨胀，小样本中几乎必然出现高相关公式。

本框架借鉴材料符号描述符发现与量纲约束思想\cite{Ouyang2018}，但把搜索空间压缩为人工审查过的规则表：同族允许加法或乘法；跨族只允许 A--B、A--D$'$、B--D$'$、A--E、A--H 等明确邻接；自动比值只允许 A/C 方向；加法要求已知维度一致；零或非有限分母保留为缺失；禁止 log、sqrt 和 power。三元式只能由同族两个代表先结合，再与一个显式相邻族代表结合。

这种规则的含义可以用实际对象解释。例如，A/C 比值可把“局域配位几何尺度”与“Na 位点数量条件”联系起来，但反向 C/A 不会因统计得分自动生成；A$\times$B 可以表示局域环境与网络连通的协同；A$+$E 只有在维度相容时才允许。每个公式保存组分、物理族、运算符、规则编号、原始值来源、缺失策略和标准化时机。搜索器无法隐藏一个公式是怎样产生的。

与 SISSO 或通用符号回归相比，这一做法可能遗漏真实但未被规则允许的关系。该损失是有意的：当前目标是获得少量可人工审查、可在结构中定位并可交给 AIMD/实验验证的假设，而不是证明所选公式是所有数学表达式中的全局最优。

\section{Stage 4：交叉验证不是一个分数，而是三个不同的科学问题}

候选公式进入预测验证后，Ridge 的角色再次改变。此时不再用于去除体系背景，而是作为低容量探针：若一个或两个描述符确实携带稳定的电导率排序信息，一个简单的线性正则模型应能在未见样本上恢复部分排序；若只有高容量非线性模型才能取得性能，则很难判断信号来自结构规律还是模型自由度。

每个训练折独立拟合中位数插补、标准化和 Ridge。测试材料的描述符分布不能提前参与中位数或尺度估计。将预处理放在全数据上再交叉验证，会使测试折信息泄漏到训练过程；在模型选择后使用同一 CV 分数作为无偏性能估计也会产生乐观偏差\cite{VarmaSimon2006}。

三种验证分别回答不同问题：

\begin{enumerate}
  \item \textbf{阴离子分层 CV：}在训练和测试具有近似化学组成比例时，候选能否对未见材料进行体系内插值？它主要检验常规预测稳定性。当前 iodide 仅一例，因此该策略必须跳过，不能生成一个看似完整的三折平均。
  \item \textbf{留一体系验证（LOSO）：}若 NASICON、sulfide 或 halide 整个体系从未进入训练，候选在该未知化学域中是否仍保持电导率排序？这是最接近本研究“跨体系可迁移”目标的测试。
  \item \textbf{重复体系分层子采样：}在 10 次 80/20 划分中，候选表现是否依赖某一次幸运划分？它估计的是划分敏感性，而不是体系外推。
\end{enumerate}

综合分只平均实际可用策略的绝对 Spearman，并同时报告可用策略数。策略跳过表示“没有证据”，不是负证据，更不能被填成零分。一个候选即使在随机子采样中表现良好，若 LOSO 方向反转，也不能称为跨体系结构规律。

\section{V1--V4：从一个相关系数扩展为多层证据}

Stage 4 的 Top 候选进一步接受四类互补审计。

\begin{enumerate}
  \item \textbf{V1 匹配公式噪声：}在每个材料体系内部打乱公式各分量，再使用完全相同的运算符重建公式。它回答“该复杂公式是否只是因自由度增加而偶然高相关”。
  \item \textbf{V2 已知因素后的折外残差预测：}训练折先根据 $Z$ 估计目标的体系背景，再让候选公式预测未见样本的 $Y$ 残差。它回答“在已知体系因素之外，公式是否仍含可折外复现的信息”。这是一种正交化预测关联，不是因果效应。
  \item \textbf{V3 体系内方向：}分别在 NASICON、sulfide 和 halide 内计算公式与 $Y$ 的 Spearman。它回答“全局关系是否完全由一个体系驱动”。
  \item \textbf{V4 体系分层 bootstrap：}在每个体系内部有放回抽样，保持体系构成并估计相关区间。它回答“固定公式的关联对材料重采样有多敏感”。
\end{enumerate}

\begin{figure}[htbp]
\centering
\resizebox{\textwidth}{!}{\begin{tikzpicture}[
  font=\sffamily\scriptsize,
  stair/.style={draw=NatureBlue!60,fill=NatureBlue!6,rounded corners=2pt,minimum height=10mm,align=left},
  arr/.style={-{Latex[length=2mm]},thick,draw=NatureGray}
]
\node[stair,minimum width=32mm] (a) at (0,0) {1. 原始相关\\$X$ 与 $Y$ 是否同向？};
\node[stair,minimum width=39mm] (b) at (1.1,1.45) {2. 去混杂\\移除体系 $Z$ 后是否保留？};
\node[stair,minimum width=46mm] (c) at (2.2,2.9) {3. 稳定性\\半样本 Lasso 是否反复选择？};
\node[stair,minimum width=53mm] (d) at (3.3,4.35) {4. 噪声与物理约束\\是否超过随机变量且公式可解释？};
\node[stair,minimum width=60mm] (e) at (4.4,5.8) {5. 跨体系验证\\LOSO 中方向和排序是否保留？};
\node[stair,minimum width=67mm] (f) at (5.5,7.25) {6. V1--V4\\随机公式、单体系驱动和不确定性是否可接受？};
\node[draw=NatureGreen!70,fill=NatureGreen!8,rounded corners=3pt,minimum width=74mm,minimum height=12mm,align=center,font=\sffamily\small\bfseries] (g) at (6.6,8.95)
{进入 AIMD / 势垒计算 / 实验验证的候选结构假设};
\draw[arr] (a.north east) -- (b.south west);
\draw[arr] (b.north east) -- (c.south west);
\draw[arr] (c.north east) -- (d.south west);
\draw[arr] (d.north east) -- (e.south west);
\draw[arr] (e.north east) -- (f.south west);
\draw[arr] (f.north east) -- (g.south west);
\node[draw=NatureRed!60,fill=NatureRed!6,rounded corners=2pt,align=center,minimum width=35mm,minimum height=18mm] at (10.4,4.2)
{任一步失败都对应\\不同的假发现来源\\而不是简单“低分”};
\end{tikzpicture}
}
\caption{\textbf{一个描述符从相关变量升级为物理验证候选所需的证据阶梯。} 任一步失败都指出一种不同的风险：体系代理、样本不稳定、随机公式、单体系驱动或跨域失效。通过全部层级仍只意味着“值得进一步验证”，并不直接证明迁移机制。}
\label{fig:evidenceladder}
\end{figure}

V1--V4 不可互相替代。一个公式可能超过置换噪声，却只在 sulfide 内成立；也可能各体系方向一致，但 bootstrap 区间很宽；还可能在固定公式条件下表现稳定，却因缺少外层“重新筛选--重新组合”的 nested reselection 而低估选择不确定性。文章必须同时呈现证据与缺口。
